\vspace{2.5pt plus 3pt minus 0.5pt}\noindent\hspace*{1.2em}\kbvnum{(6)}\hspace{0.55em}ತತ್ಪುರುಷ ಸಮಾಸಕ್ಕೂ ಕರ್ಮಧಾರಯ ಸಮಾಸಕ್ಕೂ ಇರುವ ಮುಖ್ಯ ವ್ಯತ್ಯಾಸವೇನು?\par

\vspace{2.5pt plus 3pt minus 0.5pt}\noindent\hspace*{1.2em}\kbvnum{(7)}\hspace{0.55em}ಅಂಶಿ ಸಮಾಸ ಎಂದರೇನು? ಉದಾಹರಣೆ ಸಹಿತ ವಿವರಿಸಿ.\par

\vspace{2.5pt plus 3pt minus 0.5pt}\noindent\hspace*{1.2em}\kbvnum{(8)}\hspace{0.55em}ದ್ವಿಗು ಸಮಾಸವೆಂದರೇನು? ಎರಡು ಉದಾಹರಣೆಗಳನ್ನು ಕೊಡಿ.\par

\vspace{2.5pt plus 3pt minus 0.5pt}\noindent\hspace*{1.2em}\kbvnum{(9)}\hspace{0.55em}ಬಹುವ್ರೀಹಿ ಸಮಾಸವೆಂದರೇನು? ಉದಾಹರಣೆ ಸಹಿತ ತಿಳಿಸಿ.\par

\vspace{2.5pt plus 3pt minus 0.5pt}\noindent\hspace*{1.2em}\kbvnum{(10)}\hspace{0.55em}ಕ್ರಿಯಾ ಸಮಾಸವೆಂದರೇನು? ಉದಾಹರಣೆ ಸಹಿತ ತಿಳಿಸಿ.\par

\vspace{2.5pt plus 3pt minus 0.5pt}\noindent\hspace*{1.2em}\kbvnum{(11)}\hspace{0.55em}ಗಮಕ ಸಮಾಸಕ್ಕೆ ಎರಡು ಉದಾಹರಣೆಗಳನ್ನು ಕೊಡಿ.\par

\vspace{2.5pt plus 3pt minus 0.5pt}\noindent\hspace*{1.2em}\kbvnum{(12)}\hspace{0.55em}ಇವುಗಳ ವಿಗ್ರಹವಾಕ್ಯ ಬರೆದು, ಸಮಾಸದ ಹೆಸರು ತಿಳಿಸಿ.\par

\vspace{5pt plus 4pt minus 1pt}
\noindent{\renewcommand{\arraystretch}{1.24}\setlength{\tabcolsep}{4pt}\rowcolors{1}{kbvrow}{white}
\begin{tabularx}{\textwidth}{X X X X}
ಮನೆದೇವರು & ಕೆಂದಾವರೆ & ಬಿರುಗಾಳಿ & ಹಿಂದಲೆ \\
ತಲೆದೆಗೆದು & ಚಾಕುಚೂರಿ & ಹುರಿಗಾಳು & ಕಣ್ಮುಚ್ಚಿ \\
ಕುಲದೇವತೆ & ಕಲ್ಲು & ಬಗ & ಮಳೆಬೆಳೆ \\
ಹಾಸಿಗೆಹೊದ್ದಿಕೆ & ಬಿಸಿನೀರು & ಬಯಲಾಟ & ಹೊರಮೈ \\
ಬಿಳಿದಾವರೆ & ಕಮಲಾಕ್ಷಿ & ಒಣಮರ & ಉಂಡನ್ನು \\
ಕೋದಂಡಪಾಣಿ & ಪೀತಾಂಬರ & ನಡುಮನೆ & ತಣ್ಣೀರು \\
ಕರಿತುರಗರಥ & ಕೋಳಿಮರಿ & ಒಳಗೆರೆ & ಮುಂಗೈ \\
ದಿವ್ಯಪ್ರಕಾಶ & ರೋಗಭಯ & ಮುಮ್ಮುಡಿ & ಹೆಮ್ಮಾರಿ \\
ದಾರಿಗಾಣದೆ & ಡಿಮುದು & ಕಿಬ್ಬೊಟ್ಟೆ & ಹೆದ್ದಾರಿ \\
ಮುಂಗುರುಳು & ಮೈದೋರು & ಮುಕ್ಕಣ್ಣ & ದಪ್ಪಕ್ಕೆ \\
ಮೈದಡವಿ & ಕೈಕೊಡು & ಇಮ್ಮಡಿ & ಶಶಿಧರ \\
\end{tabularx}}
\vspace{5pt plus 4pt minus 1pt}

\kbvchapter{ಪಾಠ ೧೪}{ಕೃದಂತಗಳು}

\indent ನಾಮ ಪ್ರಕೃತಿಗಳಲ್ಲಿ ಹಲವು ವಿಭೇದಗಳುಂಟು. ಇವುಗಳಲ್ಲಿ ಕೃನ್ನಾಮ ಪ್ರಕೃತಿ ಎಂಬುದು ಒಂದು ವಿಭೇದ. ಇದನ್ನು ಕೃದಂತ ನಾಮ ಪ್ರಕೃತಿ ಯೆಂತಲೂ ಕರೆಯುವುದುಂಟು.\par
